\chapter{Analisi del problema}

% ======================================================
%
% 1.1 - PROBLEM DESCRIPTION
%
% ======================================================
\section{Descrizione del problema}

% ======================================================
% 1.1.1 - N-BODY PROBLEM
% ======================================================
\subsection{Il problema N-Body}
Si vuole simulare il moto di $N$ corpi che si attraggono a vicenda per gravità, senza altre forze in gioco.
Ogni corpo ha una posizione, una velocità e una massa.
Dato lo stato iniziale, si vuole sapere come si muovono i corpi nel tempo.

Per più di due corpi non esiste una soluzione esatta, quindi il moto viene approssimato con una simulazione a passi discreti.
Nel nostro caso i corpi si muovono in uno spazio a due dimensioni e vengono generati in modo casuale.

La forza con cui il corpo $j$ attrae il corpo $i$ è
\[
  F_{ij} = G \cdot \frac{m_i \, m_j}{d_{ij}^2}
\]
dove $d_{ij}$ è la distanza tra i due corpi.
Se due corpi sono nella stessa posizione la forza sarebbe infinita, e quindi viene ignorata.

% ======================================================
% 1.1.2 - SEQUENTIAL ALGORITHM
% ======================================================
\subsection{Algoritmo sequenziale di riferimento}
L'algoritmo fornito è il seguente:
\begin{lstlisting}[basicstyle=\ttfamily\small]
vt = 0                       // tempo virtuale
dt = 0.01                    // durata di un passo
inizializza i corpi (posizione, massa, velocita')

loop {
  per ogni corpo b[i]:
    calcola la forza totale degli altri corpi b[j]
    calcola l'accelerazione (forza / massa)
    aggiorna la velocita' dopo il tempo dt
  aggiorna la posizione di ogni corpo dopo il tempo dt
  vt = vt + dt
  visualizza lo stato (se c'e' la GUI)
}
\end{lstlisting}

Ogni passo ha quindi due fasi: prima si aggiornano le velocità, poi le posizioni.
La prima è la più pesante ($O(N^2)$), ed è quella che conviene parallelizzare.

% ======================================================
%
% 1.2 - CONCURRENCY ASPECTS
%
% ======================================================
\section{Aspetti rilevanti per la concorrenza}
\label{sec:aspetti_concorrenza}

% ======================================================
% 1.2.1 - WORK DECOMPOSITION
% ======================================================
\subsection{Decomposizione del lavoro}
Dentro una fase i corpi si possono elaborare in parallelo, perché non dipendono l'uno dall'altro.
Si dividono quindi i corpi in blocchi, circa uno per core: ad esempio 5\,000 corpi su 10 core danno 500 corpi per blocco.
I passi invece vanno eseguiti uno dopo l'altro, perché ognuno parte dal risultato del precedente.

% ======================================================
% 1.2.2 - SHARED STATE
% ======================================================
\subsection{Stato condiviso}
I dati condivisi sono i corpi:
\begin{itemize}
  \item nella prima fase ogni thread legge le posizioni di \textbf{tutti} i corpi, ma scrive solo le velocità dei \textbf{propri};
  \item nella seconda fase ogni thread legge e scrive solo i \textbf{propri} corpi.
\end{itemize}
Ogni corpo viene modificato da un solo thread, quindi non serve mutua esclusione.
Bisogna però fare attenzione a \emph{quando} si scrive.

% ======================================================
% 1.2.3 - SYNCHRONIZATION BETWEEN PHASES
% ======================================================
\subsection{Sincronizzazione tra le fasi}
Servono due punti di attesa per ogni passo:
\begin{enumerate}
  \item \textbf{tra le due fasi}: nessuno può aggiornare le posizioni finché tutti non hanno finito di calcolare le velocità, altrimenti qualcuno leggerebbe posizioni già cambiate (\emph{race condition});
  \item \textbf{a fine passo}: il passo successivo e il disegno partono solo quando tutte le posizioni sono aggiornate.
\end{enumerate}
In entrambi i casi serve una \emph{barriera}: ogni thread aspetta che anche gli altri abbiano finito.

% ======================================================
%
% 1.3 - REQUIREMENTS
%
% ======================================================
\section{Requisiti}
\label{sec:requisiti}

% ======================================================
% 1.3.1 - NO-GUI VERSION
% ======================================================
\subsection{Versione senza GUI}
Esegue $M$ passi su $N$ corpi e punta solo alla velocità.
Servono due versioni: una con i \textbf{thread} e una con l'\textbf{Executor Framework}.
Per coordinare i thread si usano monitor scritti da noi.
Il risultato deve essere uguale a quello sequenziale, senza deadlock, e le prestazioni devono crescere con il numero di core.

% ======================================================
% 1.3.2 - GUI VERSION
% ======================================================
\subsection{Versione con GUI}
Mostra ogni passo e ha i comandi \textbf{start}, \textbf{stop}, \textbf{pause}, \textbf{resume} e \textbf{step} (un passo alla volta).
L'interfaccia non deve bloccarsi e deve mostrare sempre lo stato di fine passo.
